% GENERATED FILE. Edit canonical lessons, then run npm run generate:books.
% canonical-source-hash: fnv1a64:1cd3d335e17fae53
% canonical-lessons: FR-C389-enquete, FR-C389-enqueter, FR-C389-temoin, FR-C389-temoigner, FR-C389-victime, FR-R389-revision-382-389

\chapter{L'enquête}
\label{ch:fr-l-enquete}
\hlchaptermodality{389}

\begin{chapteropening}
\textbf{By the end of this chapter:} \emph{I can talk about an investigation: investigating, a witness, giving evidence and the victim.}

Complete the last lesson of chapter 389: I can talk about an investigation: investigating, a witness, giving evidence and the victim.
\end{chapteropening}

\section[l'enquête]{l'enquête — the investigation, the inquiry; the survey}

\label{lesson:FR-C389-enquete}



\begin{quote}
Before the new one: say the French for the prisoner, then the French for to punish.
\end{quote}

\subsection*{You'll want to know: l'enquête}
\textbf{l'enquête} — \textquotedblleft{}the investigation, the inquiry; the survey\textquotedblright{}.

\begin{quote}
La police a ouvert une enquête.
\end{quote}

\begin{quote}
\emph{The police have opened an investigation.}
\end{quote}

Also a survey: \emph{une enquête sur les jeunes}. Feminine.

\begin{cousinweb}
\emph{enquête} comes from Old French \emph{enqueste} (EtymDB), from Latin \emph{inquīsīta} \textquotedblleft{}things searched for\textquotedblright{}, from \emph{inquīrere} \textquotedblleft{}to search into\textquotedblright{}, built on \emph{quaerere} \textquotedblleft{}to seek\textquotedblright{} (Lewis \& Short).
\end{cousinweb}

\begin{grammarlens}[title={What you've built}]
One more word for telling, explaining and describing at B1.
\end{grammarlens}

\subsection*{Your turn}
\begin{itemize}
\raggedright
  \item \emph{Say it:} \emph{l'enquête}
  \item \emph{Say it:} \emph{l'enquête}, once more
  \item \emph{Say it:} a sentence of your own with \emph{l'enquête}
  \item \emph{Recall:} say the French for the prisoner, then the French for to punish, then say \emph{l'enquête} again
\end{itemize}

\begin{tcolorbox}[breakable,colback=teal!4,colframe=teal!35!black,title={Before you move on}]
What does l'enquête mean? (\textquotedblleft{}The investigation, the inquiry; the survey\textquotedblright{}.) Say it once more.
\end{tcolorbox}
\section[enquêter]{enquêter — to investigate}

\label{lesson:FR-C389-enqueter}



\begin{quote}
Before the new one: say the French for to punish, then the French for the investigation.
\end{quote}

\subsection*{You'll want to know: enquêter}
\textbf{enquêter} — \textquotedblleft{}to investigate\textquotedblright{}.

\begin{quote}
Deux policiers enquêtent sur l'accident.
\end{quote}

\begin{quote}
\emph{Two police officers are investigating the accident.}
\end{quote}

\emph{Enquêter sur quelque chose}.

\begin{cousinweb}
\emph{enquêter} is built on \emph{enquête} \textquotedblleft{}inquiry\textquotedblright{}, from Old French \emph{enqueste}, from Latin \emph{inquīrere} \textquotedblleft{}to search into\textquotedblright{} (Lewis \& Short), built on \emph{quaerere} \textquotedblleft{}to seek\textquotedblright{}.
\end{cousinweb}

\begin{grammarlens}[title={What you've built}]
One more word for telling, explaining and describing at B1.
\end{grammarlens}

\subsection*{Your turn}
\begin{itemize}
\raggedright
  \item \emph{Say it:} \emph{enquêter}
  \item \emph{Say it:} \emph{enquêter}, once more
  \item \emph{Say it:} a sentence of your own with \emph{enquêter}
  \item \emph{Recall:} say the French for to punish, then the French for the investigation, then say \emph{enquêter} again
\end{itemize}

\begin{tcolorbox}[breakable,colback=teal!4,colframe=teal!35!black,title={Before you move on}]
What does enquêter mean? (\textquotedblleft{}To investigate\textquotedblright{}.) Say it once more.
\end{tcolorbox}
\section[le témoin]{le témoin — the witness}

\label{lesson:FR-C389-temoin}



\begin{quote}
Before the new one: say the French for the investigation, then the French for to investigate.
\end{quote}

\subsection*{You'll want to know: le témoin}
\textbf{le témoin} — \textquotedblleft{}the witness\textquotedblright{}.

\begin{quote}
Un témoin a tout vu depuis sa fenêtre.
\end{quote}

\begin{quote}
\emph{A witness saw everything from his window.}
\end{quote}

At a wedding, \emph{le témoin} is the best man or a witness.

\begin{cousinweb}
\emph{témoin} comes from Old French \emph{tesmoing} (EtymDB), from Latin \emph{testimōnium} \textquotedblleft{}evidence, testimony\textquotedblright{}, built on \emph{testis} \textquotedblleft{}witness\textquotedblright{} (Lewis \& Short).
\end{cousinweb}

\begin{grammarlens}[title={What you've built}]
One more word for telling, explaining and describing at B1.
\end{grammarlens}

\subsection*{Your turn}
\begin{itemize}
\raggedright
  \item \emph{Say it:} \emph{le témoin}
  \item \emph{Say it:} \emph{le témoin}, once more
  \item \emph{Say it:} a sentence of your own with \emph{le témoin}
  \item \emph{Recall:} say the French for the investigation, then the French for to investigate, then say \emph{le témoin} again
\end{itemize}

\begin{tcolorbox}[breakable,colback=teal!4,colframe=teal!35!black,title={Before you move on}]
What does le témoin mean? (\textquotedblleft{}The witness\textquotedblright{}.) Say it once more.
\end{tcolorbox}
\section[témoigner]{témoigner — to testify, to give evidence}

\label{lesson:FR-C389-temoigner}



\begin{quote}
Before the new one: say the French for to investigate, then the French for the witness.
\end{quote}

\subsection*{You'll want to know: témoigner}
\textbf{témoigner} — \textquotedblleft{}to testify, to give evidence\textquotedblright{}.

\begin{quote}
Elle a accepté de témoigner au procès.
\end{quote}

\begin{quote}
\emph{She agreed to testify at the trial.}
\end{quote}

\emph{Témoigner de} — to show, to bear witness to.

\begin{cousinweb}
\emph{témoigner} comes from Old French \emph{tesmoignier} (EtymDB), formed on \emph{tesmoing} \textquotedblleft{}witness\textquotedblright{}, from Latin \emph{testimōnium} \textquotedblleft{}testimony\textquotedblright{}, built on \emph{testis} \textquotedblleft{}witness\textquotedblright{} (Lewis \& Short).
\end{cousinweb}

\begin{grammarlens}[title={What you've built}]
One more word for telling, explaining and describing at B1.
\end{grammarlens}

\subsection*{Your turn}
\begin{itemize}
\raggedright
  \item \emph{Say it:} \emph{témoigner}
  \item \emph{Say it:} \emph{témoigner}, once more
  \item \emph{Say it:} a sentence of your own with \emph{témoigner}
  \item \emph{Recall:} say the French for to investigate, then the French for the witness, then say \emph{témoigner} again
\end{itemize}

\begin{tcolorbox}[breakable,colback=teal!4,colframe=teal!35!black,title={Before you move on}]
What does témoigner mean? (\textquotedblleft{}To testify, to give evidence\textquotedblright{}.) Say it once more.
\end{tcolorbox}
\section[la victime]{la victime — the victim}

\label{lesson:FR-C389-victime}



\begin{quote}
Before the new one: say the French for the witness, then the French for to testify.
\end{quote}

\subsection*{You'll want to know: la victime}
\textbf{la victime} — \textquotedblleft{}the victim\textquotedblright{}.

\begin{quote}
La victime est encore à l'hôpital.
\end{quote}

\begin{quote}
\emph{The victim is still in hospital.}
\end{quote}

Always feminine in grammar, even for a man: \emph{il est la victime}.

\begin{cousinweb}
\emph{victime} is a learned borrowing from Latin \emph{victima} \textquotedblleft{}an animal offered in sacrifice\textquotedblright{} (EtymDB; Lewis \& Short).
\end{cousinweb}

\begin{grammarlens}[title={What you've built}]
One more word for telling, explaining and describing at B1.
\end{grammarlens}

\subsection*{Your turn}
\begin{itemize}
\raggedright
  \item \emph{Say it:} \emph{la victime}
  \item \emph{Say it:} \emph{la victime}, once more
  \item \emph{Say it:} a sentence of your own with \emph{la victime}
  \item \emph{Recall:} say the French for the witness, then the French for to testify, then say \emph{la victime} again
\end{itemize}

\begin{tcolorbox}[breakable,colback=teal!4,colframe=teal!35!black,title={Before you move on}]
What does la victime mean? (\textquotedblleft{}The victim\textquotedblright{}.) Say it once more.
\end{tcolorbox}
\section[(dialogue)]{Retour sur les chapitres 382–389}

\label{lesson:FR-R389-revision-382-389}



\begin{quote}
Say the French for to break and the offence.
\end{quote}

\subsection*{Your turn}
Say these aloud:

\begin{itemize}
\raggedright
  \item to break, the offence, to transgress, the penalty, to penalise
  \item the requirement, the criterion, eligible, the admission, to exclude
  \item to expire, acceptable, valid, to renew, the renewal
  \item the bureaucracy, the step, the procedure, the certificate, the official statement
  \item to accuse, the accusation, the thief, the complaint, the misdemeanour
  \item the trial, the defence, the verdict, guilty, innocent
  \item the sentence, to sentence, the conviction, the prisoner, to punish
  \item the investigation, to investigate, the witness, to testify, the victim
\end{itemize}

\begin{tcolorbox}[breakable,colback=teal!4,colframe=teal!35!black,title={Before you move on}]
To break? (\textbf{enfreindre}.) The offence? (\textbf{l'infraction}.) To transgress? (\textbf{transgresser}.) The penalty? (\textbf{la sanction}.) To penalise? (\textbf{sanctionner}.) The requirement? (\textbf{l'exigence}.) The criterion? (\textbf{le critère}.) Eligible? (\textbf{éligible}.) The admission? (\textbf{l'admission}.) To exclude? (\textbf{exclure}.) To expire? (\textbf{expirer}.) Acceptable? (\textbf{valable}.) Valid? (\textbf{valide}.) To renew? (\textbf{renouveler}.) The renewal? (\textbf{le renouvellement}.) The bureaucracy? (\textbf{la bureaucratie}.) The step? (\textbf{la démarche}.) The procedure? (\textbf{la procédure}.) The certificate? (\textbf{le certificat}.) The official statement? (\textbf{l'attestation}.) To accuse? (\textbf{accuser}.) The accusation? (\textbf{l'accusation}.) The thief? (\textbf{le voleur}.) The complaint? (\textbf{la plainte}.) The misdemeanour? (\textbf{le délit}.) The trial? (\textbf{le procès}.) The defence? (\textbf{la défense}.) The verdict? (\textbf{le verdict}.) Guilty? (\textbf{coupable}.) Innocent? (\textbf{innocent}.) The sentence? (\textbf{la peine}.) To sentence? (\textbf{condamner}.) The conviction? (\textbf{la condamnation}.) The prisoner? (\textbf{le détenu}.) To punish? (\textbf{punir}.) The investigation? (\textbf{l'enquête}.) To investigate? (\textbf{enquêter}.) The witness? (\textbf{le témoin}.) To testify? (\textbf{témoigner}.) The victim? (\textbf{la victime}.)
\end{tcolorbox}
